\section{Universal Algebra}
In \emph{universal algebra}, we take a birds eye view over the ideas that are common to all the specific algebraic structures (such as groups, rings, fields, etc.) that we have seen so far and many more. These common themes are: We are dealing with structures that involve \emph{sets}, \emph{operations} and \emph{axioms}. The sets constitute the backdrop, the operations define the ways in which we can combine set elements to obtain new set elements and the axioms tell us how we expect these operations to behave. For example, in the group $\mathcal{G} = (\mathbb{Z},+)$ of integers under addition, our set is the set of integers, we have one binary operation that we call addition and we have the group axioms that tell us, that this so called addition should be associative, there should be a neutral element and so on. ...TBC...


\subsection{Algebraic Structures}
An \emph{algebraic structure} (aka an algebraic system or just an algebra) consists of the following ingredients: ...TBC...


\begin{comment}

https://en.wikipedia.org/wiki/Universal_algebra
https://en.wikipedia.org/wiki/Algebraic_structure
https://en.wikipedia.org/wiki/Variety_(universal_algebra)
https://en.wikipedia.org/wiki/Operad

According to wikipedia, an algebraic structure is defined via a set and a collection of operations
on the set. Wouldn't it be more general to have a collection of sets and operations between elements
of those sets. Like, for example, the sets could be Nat and Bool and we could have operations like
the unary isEven: Nat -> Bool?

See:
https://en.wikipedia.org/wiki/Universal_algebra#Generalizations
but that also doesn't talk about collections of sets.

\end{comment}